% Complete Section 7.3, original PDF 95-96 / printed 85-86.
% block: trimmed-polynomial-forms-heading page:95 type:section
\subsection{{微分形式的修剪多项式空间}}\label{sec:trimmed-polynomial-forms}

% block: trimmed-polynomial-definition page:95 type:ordinary-paragraph
{由刚才建立的结果, 次数为 $r$ 的多项式 $k$-形式的完全空间可分解为}
\[
\begin{aligned}
\mathcal P_r\Lambda^k&=\mathcal P_{r-1}\Lambda^k\oplus\mathcal H_r\Lambda^k\\
&=\mathcal P_{r-1}\Lambda^k\oplus\kappa\mathcal H_{r-1}\Lambda^{k+1}\oplus d\mathcal H_{r+1}\Lambda^{k-1}.
\end{aligned}
\]
{若去掉最后一个直和项, 就得到介于 $\mathcal P_{r-1}\Lambda^k$ 与 $\mathcal P_r\Lambda^k$ 之间的空间, 即次数为 $r$ 的多项式形式的修剪空间:}
\[
\mathcal P_r^-\Lambda^k:=\mathcal P_{r-1}\Lambda^k+\kappa\mathcal H_{r-1}\Lambda^{k+1}.
\]
{注意, $\mathcal P_r^-\Lambda^0=\mathcal P_r\Lambda^0$, $\mathcal P_r^-\Lambda^n=\mathcal P_{r-1}\Lambda^n$, 但当 $0<k<n$ 时, $\mathcal P_r^-\Lambda^k$ 严格介于 $\mathcal P_{r-1}\Lambda^k$ 与 $\mathcal P_r\Lambda^k$ 之间. 利用 Koszul 复形的正合性和归纳法, 可计算被去掉的项 $d\mathcal H_{r+1}\Lambda^{k-1}$ 的维数 (见 [\BookNumber{11}, Theorem 3.3]). 由此得到修剪多项式空间的维数公式:}
\[
\dim\mathcal P_r^-\Lambda^k=\binom{r+n}{r+k}\binom{r+k-1}k.
\]

% block: trimmed-polynomial-dimension-comparison page:96 type:ordinary-paragraph
{与 \eqref{eq:complete-polynomial-form-dimension} 比较, 可见}
\[
\dim\mathcal P_r^-\Lambda^k=\frac r{r+k}\dim\mathcal P_r\Lambda^k.
\]
{有 $d\mathcal P_r^-\Lambda^k\subset d\mathcal P_r\Lambda^k\subset\mathcal P_{r-1}\Lambda^{k+1}\subset\mathcal P_r^-\Lambda^k$. 因而得到另一个空间经过修剪的多项式 de Rham 复形:}
\begin{equation}\label{eq:trimmed-polynomial-de-rham-complex}
\mathcal P_r^-\Lambda^0\xrightarrow{d}\mathcal P_r^-\Lambda^1\xrightarrow{d}\cdots\xrightarrow{d}\mathcal P_r^-\Lambda^n\to0.
\end{equation}

% block: trimmed-polynomial-exactness page:96 type:ordinary-paragraph
{注意, 与 Corollary~\ref{cor:polynomial-de-rham-exactness} 中的多项式 de Rham 复形不同, 在此序列中次数 $r$ 保持不变. 此复形也正合. 为说明这一点, 首先注意, 若对某个微分形式 $\omega$ 有 $d\kappa\omega=0$, 则对 $\kappa\omega$ 应用同伦公式, 得 $\kappa\omega=0$. 简言之, $d$ 在 $\kappa$ 的像上是单射. 于是}
\begin{equation}\label{eq:trimmed-polynomial-kernel}
\mathcal N(d|_{\mathcal P_{r-1}\Lambda^k})=\mathcal N(d|_{\mathcal P_r^-\Lambda^k}).
\end{equation}
{再由 $\mathcal P_r\Lambda^{k-1}=\mathcal P_r^-\Lambda^{k-1}+d\mathcal H_{r+1}\Lambda^{k-2}$, 得}
\begin{equation}\label{eq:trimmed-polynomial-range}
d\mathcal P_r\Lambda^{k-1}=d\mathcal P_r^-\Lambda^{k-1}.
\end{equation}
{由完全多项式 de Rham 复形的正合性, \eqref{eq:trimmed-polynomial-kernel} 与 \eqref{eq:trimmed-polynomial-range} 的左端相等. 因此右端也相等:}
\[
\mathcal N(d|_{\mathcal P_r^-\Lambda^k})=d\mathcal P_r^-\Lambda^{k-1}.
\]
{这就证明了 \eqref{eq:trimmed-polynomial-de-rham-complex} 的正合性.}

% block: trimmed-space-origin-independence page:96 type:remark
\Needspace{6\baselineskip}
\begin{Remark}\label{rem:trimmed-origin-independence}
{Koszul 微分定义为在每点 $x$ 处, 将微分形式与 $T_x\mathbb R^n$ 中的向量缩并; 该向量由 $\mathbb R^n\equiv T_0\mathbb R^n$ 中从原点指向 $x$ 的向量平移得到. 因此, Koszul 微分依赖于 $\mathbb R^n$ 中原点的选取. 但是可以验证 [\BookNumber{13}, p. 331], 多项式空间 $\mathcal P_r^-\Lambda^k$ 不受此选取影响. 这一观察很重要, 因为它使我们不仅能在 $\mathbb R^n$ 上, 还能在任意仿射子空间上定义修剪空间 $\mathcal P_r^-\Lambda^k$. 随后容易看出, 若 $\omega\in\mathcal P_r^-\Lambda^k(\mathbb R^n)$, 则 $\omega$ 在仿射子空间 $V$ 上的迹属于 $\mathcal P_r^-\Lambda^k(V)$.}
\end{Remark}
